Benchmark papers report scores; they rarely report whether the reported denominator can be reconstructed from the released artifacts. That gap is costly in medical vision--language model (VLM) evaluation, where one study can require hundreds of thousands of calls, hundreds of GPU-hours, and a chain of preprocessing, parsing and aggregation choices a reader cannot replay. We make the denominator the object of study. We present a harness for a three-layer chest X-ray VLM study --- in-domain NIH (1,506 studies, 19-cell degradation grid), cross-hospital PadChest (1,500) and report-generation Open-i (947), 38 stage entries, \textbf{275,218 model calls}, 39 frozen prediction files (0.49\,GiB) --- and then \emph{measure the harness against four alternatives on the same workload}: naive append, append\,+\,fsync, atomic whole-file rewrite, and SQLite in WAL mode.

Three results replace the field's usual self-report. First, \emph{cost}: at batch size 1 the shipped writer pays 1.80\,ms per cell, 4.2$\times$ a per-row fsync and 580$\times$ a naive append in p50 latency, but at the delivered batch sizes (4--64) the amortised cost is only about two minutes across the entire 275k-call study, which is negligible against inference time, and it adds no per-row storage over plain JSONL (SQLite adds 37\%). Second, \emph{fault behaviour}: after a 700-cell crash it loses 0 rows against 10 for naive append; with two writers over the same key range it fails loudly, whereas naive append and fsync-per-row silently produce 545 and 700 duplicate keys, and SQLite is the only baseline that serialises both writers cleanly; a corrupted or over-advanced sidecar costs zero rows. Third, \emph{discrimination}: our census procedure recomputes each file's denominator from the file itself, and applied to 259 prediction files it flags 4 non-uniform denominators and one file with duplicate keys, so the audit is capable of failing.

A remote arm of 9,400 calls against a frontier-scale open-weight model, served through an OpenAI-compatible relay, adds three provider-side failure classes --- hidden reasoning exhausting the token budget, TLS resets under concurrency, and a billing wall --- all detected by the same contract; after recovery the file carries 9,400 unique keys, one provenance id and zero transport errors.

We replace the unfalsifiable claim ``zero error rows'' with five testable invariants --- unique keys, no acknowledged-row loss, no silent I/O success, no torn tail accepted as complete, no silent configuration mixture --- and report which of the five baselines satisfies each. The result is not that our harness is durable in the abstract, but that durability has a measured price and a measured payoff, and that a benchmark's denominator can be made checkable by a third party.
